\documentclass[12pt]{article}
\usepackage[utf8]{inputenc}
\usepackage[margin=1in]{geometry}
\usepackage{amsmath, amssymb, amsthm}
\usepackage{xcolor}
\usepackage{tcolorbox}
\tcbuselibrary{breakable}
\usepackage{enumitem}
\usepackage{fancyhdr}
\usepackage{titlesec}
\usepackage{hyperref}

% Define colored boxes - all with breakable option
\newtcolorbox{configbox}[1]{
    colback=blue!5!white,
    colframe=blue!75!black,
    title=#1,
    fonttitle=\bfseries,
    breakable
}

\newtcolorbox{strategybox}[1]{
    colback=pink!10!white,
    colframe=pink!75!black,
    title=#1,
    fonttitle=\bfseries,
    breakable
}

\newtcolorbox{tacticsbox}[1]{
    colback=green!10!white,
    colframe=green!75!black,
    title=#1,
    fonttitle=\bfseries,
    breakable
}

\newtcolorbox{conversionbox}[1]{
    colback=yellow!10!white,
    colframe=yellow!75!black,
    title=#1,
    fonttitle=\bfseries,
    breakable
}

\newtcolorbox{verificationbox}[1]{
    colback=red!10!white,
    colframe=red!75!black,
    title=#1,
    fonttitle=\bfseries,
    breakable
}

\newtcolorbox{roadmapbox}[1]{
    colback=cyan!10!white,
    colframe=cyan!75!black,
    title=#1,
    fonttitle=\bfseries,
    breakable
}

\newtcolorbox{competitionbox}[1]{
    colback=violet!10!white,
    colframe=violet!75!black,
    title=#1,
    fonttitle=\bfseries,
    breakable
}

\newtcolorbox{highlightbox}[1]{
    colback=orange!10!white,
    colframe=orange!75!black,
    title=#1,
    fonttitle=\bfseries,
    breakable
}

\newtcolorbox{warningbox}[1]{
    colback=red!15!white,
    colframe=red!85!black,
    title=#1,
    fonttitle=\bfseries,
    breakable
}

\newtcolorbox{protipbox}[1]{
    colback=purple!10!white,
    colframe=purple!75!black,
    title=#1,
    fonttitle=\bfseries,
    breakable
}

\title{\textbf{AMC 12A 2024 Problem 25:\\Comprehensive Strategic Analysis}}
\author{Using Enhanced Framework v2.1}
\date{\today}

\begin{document}

\maketitle

\tableofcontents
\newpage

\section{Problem Statement}

\textbf{2024 AMC 12A Problem 25}

A graph is \textit{symmetric} about a line if the graph remains unchanged after reflection in that line. For how many quadruples of integers $(a,b,c,d)$, where $|a|,|b|,|c|,|d| \leq 5$ and $c$ and $d$ are not both $0$, is the graph of
\[y = \frac{ax+b}{cx+d}\]
symmetric about the line $y=x$?

\[\textbf{(A) }1282\qquad\textbf{(B) }1292\qquad\textbf{(C) }1310\qquad\textbf{(D) }1320\qquad\textbf{(E) }1330\]

\vspace{0.5cm}

\section{Problem Configuration Analysis}

\begin{configbox}{A. Problem Classification}
\textbf{Problem Type:} To Find (counting problem)

\textbf{Difficulty Band:} Late-band (P25/25) - highest difficulty

\textbf{Competition Context:} AMC12 (75 minutes, 25 problems, multiple choice)

\textbf{Mathematical Domain:} 
\begin{itemize}
    \item Primary: Algebra (rational functions, functional equations)
    \item Secondary: Combinatorics (counting with constraints)
    \item Tertiary: Geometry (symmetry transformations)
\end{itemize}

\textbf{Source Context:} Final problem of AMC 12A, expected solve rate $<$5\%, typical time allocation 6-8 minutes
\end{configbox}

\begin{configbox}{B. Problem Structure Identification}
\textbf{Given Information:}
\begin{itemize}
    \item Rational function $y = \frac{ax+b}{cx+d}$ with integer coefficients
    \item Constraint: $|a|, |b|, |c|, |d| \leq 5$ (11 choices each: $-5$ to $5$)
    \item Constraint: $c$ and $d$ are not both $0$ (denominator must be non-constant)
    \item Symmetry condition: graph symmetric about $y=x$
\end{itemize}

\textbf{Constraints:}
\begin{itemize}
    \item Integer coefficients only
    \item Bounded magnitude: $|a|, |b|, |c|, |d| \leq 5$
    \item Non-trivial denominator: $(c,d) \neq (0,0)$
    \item Function must be well-defined (no division by zero everywhere)
\end{itemize}

\textbf{Objective:} Count all valid quadruples $(a,b,c,d)$ satisfying the symmetry condition

\textbf{Hidden Assumptions:}
\begin{itemize}
    \item Symmetry about $y=x$ is equivalent to $f(f(x)) = x$ (inverse function equals original)
    \item Must exclude degenerate cases (undefined functions, trivial constants)
    \item The function notation implicitly requires proper domain
\end{itemize}

\textbf{Answer Choice Leverage:}
\begin{itemize}
    \item All choices near 1300, suggesting systematic counting with small corrections
    \item Difference between choices is 10-30, indicating edge cases matter significantly
    \item Can verify final answer against choices to catch arithmetic errors
\end{itemize}
\end{configbox}

\begin{configbox}{C. Strategic Orientation}
\textbf{Hypothesis:} A rational function $y = \frac{ax+b}{cx+d}$ with integer coefficients bounded by 5

\textbf{Conclusion:} The count of valid quadruples satisfying symmetry about $y=x$

\textbf{Notation:}
\begin{itemize}
    \item $f(x) = \frac{ax+b}{cx+d}$ (the given function)
    \item $f^{-1}(x)$ (inverse function)
    \item Symmetry condition: $f(f(x)) = x$ for all $x$ in domain
\end{itemize}

\textbf{Intuitive Approach:} 
\begin{itemize}
    \item Symmetry about $y=x$ means reflecting over the line $y=x$ gives the same graph
    \item Equivalently, swapping $x$ and $y$ in the equation gives the same relation
    \item This means $f^{-1}(x) = f(x)$, so the function is its own inverse
\end{itemize}

\textbf{Cross-Domain Potential:}
\begin{itemize}
    \item Geometric approach: Use transformation properties (rotation by $45^\circ$)
    \item Algebraic approach: Derive conditions from $f(f(x)) = x$
    \item Asymptotic approach: Use properties of hyperbola asymptotes
    \item Combinatorial approach: Case analysis with systematic counting
\end{itemize}
\end{configbox}

\section{Strategic Framework Application}

\begin{strategybox}{A. Psychological Strategies}
\textbf{Mental Toughness:}
\begin{itemize}
    \item This is P25 - expect multiple failed attempts
    \item Complex case analysis required - stay organized
    \item Must handle both algebraic manipulation and combinatorial counting
    \item Don't get discouraged by messy algebra in early attempts
\end{itemize}

\textbf{Creativity Cultivation:}
\begin{itemize}
    \item Consider multiple interpretations of "symmetry about $y=x$"
    \item Look for shortcuts using answer choices
    \item Think about special cases: $c=0$ (linear), $a=-d$ (hyperbola)
    \item Use geometric intuition to guide algebraic work
\end{itemize}

\textbf{Iterative Refinement:}
\begin{itemize}
    \item Start with understanding $f^{-1} = f$ condition
    \item Refine by splitting into cases based on $c$
    \item Further refine by handling degenerate cases
    \item Verify counts against answer choices
\end{itemize}
\end{strategybox}

\begin{strategybox}{B. Investigation Strategies}
\textbf{Startup Orientation:}
\begin{itemize}
    \item Recognize this as a functional equation problem disguised as counting
    \item Key insight: Symmetry about $y=x$ $\Leftrightarrow$ $f$ is self-inverse
    \item Main task: Find algebraic conditions, then count systematically
\end{itemize}

\textbf{Get Your Hands Dirty:}
\begin{itemize}
    \item Test simple cases: $y=x$ (trivial symmetry), $y=-x+k$ (perpendicular lines)
    \item Try $c=0$: linear functions, easier to analyze
    \item Try $c \neq 0$: general rational functions, use inverse formula
\end{itemize}

\textbf{Penultimate Step:}
\begin{itemize}
    \item Before final count: identify all valid cases
    \item Before answer: subtract overcounted/invalid cases
    \item Final arithmetic check against answer choices
\end{itemize}

\textbf{Wishful Thinking:}
\begin{itemize}
    \item Wish: condition simplifies to nice algebraic relation
    \item Hope: most cases follow simple pattern ($a=-d$)
    \item Desire: edge cases are easy to identify and count
\end{itemize}

\textbf{Make It Easier:}
\begin{itemize}
    \item Simpler problem: What if $c=0$? (linear functions only)
    \item Simpler problem: What if no bounds on coefficients?
    \item Simpler problem: What conditions for symmetry about $y=0$?
\end{itemize}

\textbf{Conjecture via Small Cases:}
\begin{itemize}
    \item Test $a=1, b=0, c=1, d=-1$: gives $y=\frac{x}{x-1}$, is it symmetric?
    \item Compute inverse: swap $x$ and $y$, solve for $y$
    \item Verify pattern: $a=-d$ seems to work for many cases
\end{itemize}

\textbf{Visualization:}
\begin{itemize}
    \item Sketch line $y=x$ and a few hyperbolas
    \item Observe: asymptotes of symmetric hyperbola intersect on $y=x$
    \item For linear case: only $y=x$ and $y=-x+k$ are symmetric
\end{itemize}
\end{strategybox}

\begin{strategybox}{C. Late-Band Strategic Playbook (P25 of 25)}
\textbf{Primary Strategy: Algebraic Decomposition + Systematic Counting}
\begin{itemize}
    \item This problem requires finding conditions first, then counting
    \item Cannot rely solely on answer-choice testing
    \item Must develop complete case analysis
\end{itemize}

\textbf{Answer-Choice Leverage:}
\begin{itemize}
    \item Use to verify final arithmetic
    \item All answers near 1300 $\Rightarrow$ main count is $11^3 = 1331$
    \item Small corrections $\Rightarrow$ edge cases subtract 30-50 invalid cases
\end{itemize}

\textbf{Make It Easier:}
\begin{itemize}
    \item Start with $c=0$ case (linear functions)
    \item This gives clean conditions: $a=\pm d$ with restrictions on $b$
    \item Then tackle $c \neq 0$ case with more general theory
\end{itemize}

\textbf{Penultimate Step:}
\begin{itemize}
    \item Main condition: $a=-d$ (works for vast majority)
    \item Count: $11^3 = 1331$ if only this condition
    \item Subtract: invalid cases where function is not well-defined or not symmetric
\end{itemize}
\end{strategybox}

\begin{strategybox}{D. Argument Construction Strategy}
\textbf{Direct Proof Approach:}
\begin{itemize}
    \item Show: Symmetry about $y=x$ $\Leftrightarrow$ $f^{-1}(x) = f(x)$
    \item Derive: Necessary conditions on $(a,b,c,d)$
    \item Verify: These conditions are sufficient
    \item Count: Valid quadruples satisfying all conditions and constraints
\end{itemize}

\textbf{Case Analysis:}
\begin{itemize}
    \item \textbf{Case 1:} $c=0$ (linear functions)
    \begin{itemize}
        \item Subcase 1.1: $b \neq 0$, must have $a=-d$
        \item Subcase 1.2: $b=0$, must have $a^2 = d^2$
    \end{itemize}
    \item \textbf{Case 2:} $c \neq 0$ (general rational functions)
    \begin{itemize}
        \item Main condition: $a=-d$
        \item Exclusions: $bc-ad=0$ (reduces to constant function)
    \end{itemize}
\end{itemize}
\end{strategybox}

\section{Tactical Methodology Selection}

\begin{tacticsbox}{A. Core Mathematical Tactics}
\textbf{Relevant Tactics for This Problem:}

\textbf{1. Symmetry Exploitation:}
\begin{itemize}
    \item Symmetry about $y=x$ is reflection symmetry
    \item Equivalent to function being its own inverse
    \item Use: $f(f(x)) = x$ as working condition
\end{itemize}

\textbf{2. Invariant Recognition:}
\begin{itemize}
    \item The property $f^{-1} = f$ is the key invariant
    \item Under reflection about $y=x$: $(x,y) \mapsto (y,x)$
    \item This swaps the roles of independent and dependent variables
\end{itemize}

\textbf{3. Extreme Principle:}
\begin{itemize}
    \item Test extreme cases: $a=b=c=d=0$ (excluded), $a=\pm 5$, etc.
    \item Boundary behavior: when does function become undefined?
    \item Limiting cases: linear functions ($c=0$), constant functions
\end{itemize}

\textbf{4. Wishful Thinking:}
\begin{itemize}
    \item Wish for simple condition like $a=-d$
    \item This turns out to be the main condition!
    \item Reduces counting to manageable cases
\end{itemize}

\textbf{5. Case Analysis:}
\begin{itemize}
    \item Split on whether $c=0$ or $c \neq 0$
    \item Further split on special values: $b=0$, $d=0$, etc.
    \item Systematic counting within each case
\end{itemize}
\end{tacticsbox}

\begin{tacticsbox}{B. Domain-Specific Techniques}
\textbf{Algebra Techniques:}
\begin{itemize}
    \item Finding inverse of rational function
    \item Solving functional equation $f(f(x)) = x$
    \item Polynomial identity (equating coefficients)
    \item Algebraic manipulation to isolate conditions
\end{itemize}

\textbf{Combinatorics Techniques:}
\begin{itemize}
    \item Systematic case enumeration
    \item Inclusion-exclusion to handle constraints
    \item Counting with restrictions (e.g., $c$ and $d$ not both 0)
    \item Accounting for symmetry in counting (e.g., $a=\pm d$)
\end{itemize}

\textbf{Geometry Techniques:}
\begin{itemize}
    \item Understanding reflection about line $y=x$
    \item Asymptote analysis for hyperbolas
    \item Coordinate transformation (rotation by $45^\circ$)
\end{itemize}
\end{tacticsbox}

\begin{tacticsbox}{C. Cross-Domain Connections}
\textbf{Algebra $\leftrightarrow$ Geometry:}
\begin{itemize}
    \item Algebraic condition $a=-d$ $\Leftrightarrow$ Geometric property that asymptotes intersect on $y=x$
    \item Functional equation $\Leftrightarrow$ Symmetry property
\end{itemize}

\textbf{Algebra $\leftrightarrow$ Combinatorics:}
\begin{itemize}
    \item Algebraic conditions determine valid cases
    \item Combinatorial counting to find total number
\end{itemize}

\textbf{Multiple Solution Paths:}
\begin{itemize}
    \item \textbf{Path 1 (Inverse Function):} Derive $f^{-1}$, set equal to $f$
    \item \textbf{Path 2 (Rotation):} Rotate by $45^\circ$, use even function property
    \item \textbf{Path 3 (Asymptotes):} Use asymptote intersection property
    \item \textbf{Path 4 (Composition):} Expand $f(f(x))=x$ directly
\end{itemize}
\end{tacticsbox}

\section{Conversion Techniques Application}

\begin{conversionbox}{A. Recognition Index}
\textbf{Pattern Recognition:}
\begin{itemize}
    \item \textbf{Signature:} Rational function with symmetry condition
    \item \textbf{Key Feature:} "Symmetric about $y=x$"
    \item \textbf{Trigger:} Self-inverse function, $f^{-1}=f$
    \item \textbf{Similar Problems:} Involutions, fixed points, functional equations
\end{itemize}

\textbf{Domain Identification:}
\begin{itemize}
    \item Primary domain: Algebra (rational functions)
    \item Natural conversion: Geometry (reflections)
    \item Alternative conversion: Functional analysis (involutions)
    \item Counting aspect: Combinatorics
\end{itemize}
\end{conversionbox}

\begin{conversionbox}{B. Technique Selection}
\textbf{Primary Approach: Inverse Function Method}

\textbf{Why this approach:}
\begin{itemize}
    \item Most direct algebraic path
    \item Clearly separates cases ($c=0$ vs $c \neq 0$)
    \item Leads to explicit conditions on coefficients
    \item Standard technique for AMC-level problems
\end{itemize}

\textbf{Key steps:}
\begin{enumerate}
    \item Find $f^{-1}(x)$ by swapping $x$ and $y$, solving for $y$
    \item Set $f^{-1}(x) = f(x)$ and simplify
    \item Extract conditions on $a,b,c,d$
    \item Count valid quadruples with systematic case analysis
\end{enumerate}

\textbf{Alternative Approaches:}
\begin{itemize}
    \item \textbf{Rotation Method:} More elegant but requires complex number machinery
    \item \textbf{Asymptote Method:} Good geometric intuition but requires careful case work
    \item \textbf{Direct Composition:} Expand $f(f(x))=x$ but gets algebraically messy
\end{itemize}
\end{conversionbox}

\begin{conversionbox}{C. Integration Strategy}
\textbf{Combining Multiple Techniques:}

\textbf{Phase 1: Algebraic Derivation}
\begin{itemize}
    \item Use inverse function method to get conditions
    \item Result: Two main cases based on $c$
\end{itemize}

\textbf{Phase 2: Combinatorial Counting}
\begin{itemize}
    \item For each case, count systematically
    \item Use multiplication principle for independent choices
    \item Subtract invalid/degenerate cases
\end{itemize}

\textbf{Phase 3: Geometric Verification}
\begin{itemize}
    \item Check answer against geometric intuition
    \item Verify edge cases make sense visually
    \item Use asymptote reasoning as sanity check
\end{itemize}
\end{conversionbox}

\begin{conversionbox}{D. Conflict Resolution}
\textbf{Potential Conflicts:}

\textbf{Issue 1: Degenerate cases}
\begin{itemize}
    \item Problem: When is function not well-defined?
    \item Resolution: Identify when $cx+d$ is identically zero or when numerator/denominator have common factor
    \item Action: Subtract these from main count
\end{itemize}

\textbf{Issue 2: Overcounting}
\begin{itemize}
    \item Problem: Some conditions might count same case multiple ways
    \item Resolution: Ensure cases are mutually exclusive
    \item Action: Carefully structure case divisions
\end{itemize}

\textbf{Issue 3: Boundary conditions}
\begin{itemize}
    \item Problem: Edge cases at coefficient bounds ($\pm 5$)
    \item Resolution: Check that counting includes/excludes these correctly
    \item Action: Verify count formulas at extremes
\end{itemize}
\end{conversionbox}

\section{Verification Strategy Design}

\begin{verificationbox}{A. Verification Level Selection}
\textbf{Recommended Level: Level 4-5 (8-12 minutes total problem time)}

\textbf{Decision Tree:}
\begin{itemize}
    \item Problem 25 (highest difficulty) $\Rightarrow$ High stakes
    \item Counting problem $\Rightarrow$ Arithmetic errors likely
    \item Multiple cases $\Rightarrow$ Need systematic verification
    \item Multiple choice $\Rightarrow$ Can check against answers
\end{itemize}

\textbf{Recommended Verification:}
\begin{itemize}
    \item \textbf{Level 1 (30s):} Check answer is among choices
    \item \textbf{Level 3 (2-3 min):} Verify main cases arithmetically
    \item \textbf{Level 4 (5-6 min):} Check a specific example from each case
    \item \textbf{Level 5 (8-10 min):} Alternative method verification (if time permits)
\end{itemize}
\end{verificationbox}

\begin{verificationbox}{B. Specialized Verification Methods}
\textbf{For This Problem:}

\textbf{1. Example Verification:}
\begin{itemize}
    \item Pick $(a,b,c,d) = (1,2,3,-1)$ from Case 2
    \item Compute: $f(x) = \frac{x+2}{3x-1}$
    \item Find inverse: $f^{-1}(x) = \frac{x+2}{3x-1}$ (swap and solve)
    \item Verify: $f^{-1}(x) = f(x)$ $\checkmark$
\end{itemize}

\textbf{2. Extreme Case Verification:}
\begin{itemize}
    \item Check $c=0, b=5, d=-5, a=5$: $y = -x+1$ (symmetric) $\checkmark$
    \item Check $c=0, b=0, d=5, a=5$: $y = x$ (symmetric) $\checkmark$
    \item Check $c=0, b=0, d=5, a=-5$: $y = -x$ (symmetric) $\checkmark$
\end{itemize}

\textbf{3. Counting Verification:}
\begin{itemize}
    \item Case 1.1: $10 \times 10 = 100$ $\checkmark$ (verified)
    \item Case 1.2: $10 \times 2 = 20$ $\checkmark$ (verified)
    \item Case 2: $1331 - 28 - 11 - 10 - 10 = 1172$ (check arithmetic)
    \item Total: $100 + 20 + 1172 = 1292$ $\checkmark$ matches (B)
\end{itemize}

\textbf{4. Asymptotic Verification (Geometry):}
\begin{itemize}
    \item For hyperbola to be symmetric about $y=x$:
    \item Horizontal asymptote: $y = a/c$
    \item Vertical asymptote: $x = -d/c$
    \item Must have: $a/c = -d/c \Rightarrow a = -d$ $\checkmark$
\end{itemize}

\textbf{5. Answer Choice Elimination:}
\begin{itemize}
    \item Too low (A: 1282): Missing some valid cases
    \item Just right (B: 1292): Matches our calculation
    \item Too high (C,D,E): Overcounting invalid cases
\end{itemize}
\end{verificationbox}

\begin{verificationbox}{C. Confidence Calibration}
\textbf{Confidence Assessment:}

\textbf{High Confidence Indicators (8/10):}
\begin{itemize}
    \item Main condition $a=-d$ is well-established
    \item Multiple solution paths converge to same answer
    \item Answer matches choice (B)
    \item Arithmetic verified multiple times
    \item Edge cases systematically identified
\end{itemize}

\textbf{Remaining Uncertainties:}
\begin{itemize}
    \item Possible arithmetic error in Case 2 subcounting
    \item Could have missed a subtle degenerate case
    \item Time pressure might cause rushed counting
\end{itemize}

\textbf{Final Decision:}
\begin{itemize}
    \item Mark answer: \textbf{(B) 1292}
    \item Confidence: 80-85\%
    \item If time permits: double-check Case 2 arithmetic
\end{itemize}
\end{verificationbox}

\section{Solution Roadmap}

\begin{roadmapbox}{Step-by-Step Solution Plan}

\subsection*{Phase 1: Understand the Condition (2 minutes)}

\textbf{Step 1: Interpret "symmetric about $y=x$"}
\begin{itemize}
    \item Reflecting the graph over line $y=x$ gives the same graph
    \item Equivalently: swapping $x$ and $y$ in equation gives same relation
    \item Algebraically: $f(f(x)) = x$ or $f^{-1}(x) = f(x)$
\end{itemize}

\textbf{Checkpoint 1:} Do I understand what symmetry about $y=x$ means algebraically?

\subsection*{Phase 2: Derive Main Conditions (3-4 minutes)}

\textbf{Step 2: Find inverse function}

For $y = \frac{ax+b}{cx+d}$, swap $x$ and $y$:
\[x = \frac{ay+b}{cy+d}\]

Solve for $y$:
\begin{align*}
x(cy+d) &= ay+b\\
cxy + dx &= ay + b\\
cxy - ay &= b - dx\\
y(cx-a) &= b-dx\\
y &= \frac{b-dx}{cx-a} = \frac{-dx+b}{cx-a}
\end{align*}

So $f^{-1}(x) = \frac{-dx+b}{cx-a}$

\textbf{Step 3: Set $f^{-1}(x) = f(x)$}
\[\frac{-dx+b}{cx-a} = \frac{ax+b}{cx+d}\]

\textbf{Step 4: Cross-multiply and equate coefficients}

\textbf{Checkpoint 2:} Have I correctly derived the inverse and set up the equality?

\subsection*{Phase 3: Case Analysis (4-5 minutes)}

\textbf{Case 1: $c=0$ (Linear Functions)}

When $c=0$: $y = \frac{a}{d}x + \frac{b}{d}$

Inverse: $y = \frac{d}{a}x - \frac{b}{a}$

For these to be equal:
\begin{itemize}
    \item Slopes equal: $\frac{a}{d} = \frac{d}{a} \Rightarrow a^2 = d^2 \Rightarrow a = \pm d$
    \item Intercepts equal: $\frac{b}{d} = -\frac{b}{a}$
\end{itemize}

From intercept condition: $ab = -bd \Rightarrow b(a+d) = 0$

\textbf{Subcase 1.1:} $b \neq 0 \Rightarrow a+d=0 \Rightarrow a=-d$
\begin{itemize}
    \item Choices: 10 values for $b$ (non-zero), 10 for $d$ (non-zero, determines $a$)
    \item Count: $10 \times 10 = 100$
\end{itemize}

\textbf{Subcase 1.2:} $b=0 \Rightarrow a^2=d^2$
\begin{itemize}
    \item If $a=d$: Function is $y=x$, need $d \neq 0$
    \item If $a=-d$: Already counted in Subcase 1.1 when $b=0$... wait, no!
    \item Actually $b=0$ is separate: $a=d$ or $a=-d$, both with $b=0$, $c=0$, $d \neq 0$
    \item For each of 10 non-zero values of $d$: 2 choices for $a$ (i.e., $\pm d$)
    \item But $a=-d$ with $b=0$ gives slope $-1$, not same as $a=d$
    \item Count: $10 \times 2 = 20$
\end{itemize}

\textbf{Case 1 Total:} $100 + 20 = 120$

\textbf{Case 2: $c \neq 0$ (General Rational Functions)}

From the equality $\frac{-dx+b}{cx-a} = \frac{ax+b}{cx+d}$:

Rewrite both in form $\text{constant} + \frac{\text{remainder}}{\text{denominator}}$:
\begin{align*}
f(x) &= \frac{a}{c} + \frac{b - \frac{ad}{c}}{cx+d}\\
f^{-1}(x) &= -\frac{d}{c} + \frac{b - \frac{ad}{c}}{cx-a}
\end{align*}

For these to be equal for all $x$:
\begin{itemize}
    \item Constants equal: $\frac{a}{c} = -\frac{d}{c} \Rightarrow a = -d$
    \item Denominators equal: $cx+d = cx-a \Rightarrow d = -a$ $\checkmark$ (consistent)
\end{itemize}

Main condition: $a=-d$ with $c \neq 0$

Count without restrictions: $11 \times 11 \times 10 = 1210$
\begin{itemize}
    \item 11 choices for $a$ (including 0)
    \item 11 choices for $b$ (including 0)  
    \item 10 choices for $c$ (excluding 0, determines $d=-a$)
\end{itemize}

\textbf{But we must exclude invalid cases!}

\textbf{Exclusion 1: Function reduces to constant}

If $bc - ad = 0$ with $a=-d$, then $bc + a^2 = 0$, so numerator and denominator have common factor.

When does $bc = -a^2$?

\begin{itemize}
    \item $a=0$: $bc=0$, so $b=0$ or $c=0$. But $c \neq 0$ in Case 2, so need $b=0$. Count: 10
    \item $a=\pm 1$: $bc=-1$. Pairs: $(b,c) \in \{(1,-1), (-1,1)\}$. Count: $2 \times 2 \times 1 = 4$
    \item $a=\pm 2$: $bc=-4$. Pairs: $(1,-4), (-1,4), (2,-2), (-2,2), (4,-1), (-4,1)$. Count: $2 \times 6 = 12$
    \item $a=\pm 3$: $bc=-9$. Pairs: $(3,-3), (-3,3)$. Count: $2 \times 2 = 4$
    \item $a=\pm 4$: $bc=-16$. Pairs: $(4,-4), (-4,4)$. Count: $2 \times 2 = 4$
    \item $a=\pm 5$: $bc=-25$. Pairs: $(5,-5), (-5,5)$. Count: $2 \times 2 = 4$
\end{itemize}

Total exclusions: $10 + 4 + 12 + 4 + 4 + 4 = 38$

\textbf{Case 2 Total:} $1210 - 38 = 1172$

\textbf{Checkpoint 3:} Have I correctly counted all cases and exclusions?

\subsection*{Phase 4: Final Count and Verification (1-2 minutes)}

\textbf{Step 5: Add all cases}
\begin{align*}
\text{Total} &= \text{Case 1} + \text{Case 2}\\
&= 120 + 1172\\
&= 1292
\end{align*}

\textbf{Step 6: Verify against answer choices}

Answer is \textbf{(B) 1292} $\checkmark$

\textbf{Checkpoint 4:} Does this match an answer choice? Have I double-checked arithmetic?

\end{roadmapbox}

\begin{highlightbox}{Key Insights}
\textbf{Critical Observations:}
\begin{enumerate}
    \item Symmetry about $y=x$ means $f^{-1}(x) = f(x)$ (self-inverse)
    \item Main algebraic condition: $a=-d$ (works for vast majority of cases)
    \item Must split into $c=0$ and $c \neq 0$ for systematic counting
    \item Many degenerate cases must be excluded (when function reduces to constant)
    \item Counting requires careful casework on when $bc = -a^2$
\end{enumerate}
\end{highlightbox}

\begin{warningbox}{Common Pitfalls}
\textbf{Danger Zones:}
\begin{enumerate}
    \item \textbf{Forgetting $c=0$ case:} Linear functions are special!
    \item \textbf{Not excluding degenerate cases:} When $bc-ad=0$, function is not well-defined
    \item \textbf{Arithmetic errors:} Lots of counting, easy to make mistakes
    \item \textbf{Overcounting $b=0$ case:} Be careful with Subcase 1.2
    \item \textbf{Missing factorization pairs:} For $bc=-a^2$, systematically find all $(b,c)$ pairs
\end{enumerate}
\end{warningbox}

\begin{protipbox}{Pro Tips}
\textbf{Time-Saving Strategies:}
\begin{enumerate}
    \item Recognize self-inverse immediately from "symmetric about $y=x$"
    \item Use answer choices to guide counting (all near 1300)
    \item Main count is $11^3 = 1331$ with corrections
    \item For Case 2, use multiplication principle systematically
    \item Double-check factorizations of $-a^2$ for each $a$
\end{enumerate}
\end{protipbox}

\section{Competition Strategy Integration}

\begin{competitionbox}{A. Time Management}
\textbf{Problem 25 Time Allocation (Total: 6-8 minutes)}

\textbf{Phase 1: Initial Analysis (1 min)}
\begin{itemize}
    \item Read and understand problem
    \item Recognize self-inverse condition
    \item Plan case split
\end{itemize}

\textbf{Phase 2: Algebraic Work (2-3 min)}
\begin{itemize}
    \item Derive inverse function
    \item Extract conditions ($a=-d$, etc.)
    \item Identify main cases
\end{itemize}

\textbf{Phase 3: Counting (2-3 min)}
\begin{itemize}
    \item Count Case 1 subcases
    \item Count Case 2 with exclusions
    \item Add up total
\end{itemize}

\textbf{Phase 4: Verification (1-2 min)}
\begin{itemize}
    \item Check arithmetic
    \item Verify against answer choices
    \item Test one example if time
\end{itemize}

\textbf{Decision Point (at 4 minutes):}
\begin{itemize}
    \item If stuck on Case 2 exclusions: estimate and move on
    \item If completely stuck: skip and return if time at end
    \item If progressing well: continue to completion
\end{itemize}
\end{competitionbox}

\begin{competitionbox}{B. Strategic Prioritization}
\textbf{Problem 25 Context:}
\begin{itemize}
    \item This is the hardest problem (P25/25)
    \item Only attempt if you've finished P1-24 or nearly finished
    \item Expected solve rate: 3-5\% of test-takers
    \item Worth same 1 point as easier problems
\end{itemize}

\textbf{When to Attempt:}
\begin{itemize}
    \item $\checkmark$ If you have 10+ minutes remaining
    \item $\checkmark$ If you're comfortable with functional equations
    \item $\checkmark$ If you're strong at systematic counting
    \item $\times$ If you have fewer than 5 minutes
    \item $\times$ If you haven't finished P1-22
\end{itemize}

\textbf{Risk Assessment:}
\begin{itemize}
    \item \textbf{High reward:} Demonstrates mastery on hardest problem
    \item \textbf{High risk:} Could waste 8 minutes for no points
    \item \textbf{Alternative:} Double-check P20-24 instead
\end{itemize}
\end{competitionbox}

\begin{competitionbox}{C. Answer Recording and Format}
\textbf{Recording Strategy:}
\begin{itemize}
    \item Bubble answer immediately after verification
    \item Mark problem for review if time permits
    \item Note: $(a,b,c,d)$ are integers, answer is integer count
    \item Double-check you're answering the right question (count of quadruples, not specific quadruple)
\end{itemize}

\textbf{Format Requirements:}
\begin{itemize}
    \item AMC12: Multiple choice, bubble one answer
    \item Answer: \textbf{(B) 1292}
    \item Units: Count (dimensionless)
\end{itemize}
\end{competitionbox}

\begin{competitionbox}{D. Abandonment Criteria}
\textbf{When to Give Up:}
\begin{enumerate}
    \item \textbf{Time threshold:} If you've spent 10+ minutes without progress
    \item \textbf{Confidence threshold:} If you can't verify any cases
    \item \textbf{Progress threshold:} If stuck on deriving main condition after 3-4 minutes
\end{enumerate}

\textbf{Strategic Retreat:}
\begin{itemize}
    \item Don't be discouraged - this is expected to be very hard
    \item Better to move on and double-check earlier problems
    \item Can return if time at end (but prioritize reviewing P20-24 first)
\end{itemize}

\textbf{Psychological Management:}
\begin{itemize}
    \item Skipping P25 is normal and smart time management
    \item Focus on maximizing overall score, not solving every problem
    \item Even top students often skip P25
\end{itemize}
\end{competitionbox}

\section{Detailed Solution}

\subsection{Method 1: Inverse Function Approach}

\textbf{Step 1: Establish the equivalence}

A function is symmetric about $y=x$ if and only if it equals its inverse function:
\[f^{-1}(x) = f(x)\]

\textbf{Step 2: Find the inverse}

For $y = \frac{ax+b}{cx+d}$, swap $x$ and $y$ and solve:
\begin{align*}
x &= \frac{ay+b}{cy+d}\\
x(cy+d) &= ay+b\\
cxy + dx &= ay + b\\
y(cx-a) &= b-dx\\
y &= \frac{b-dx}{cx-a}
\end{align*}

Therefore: $f^{-1}(x) = \frac{-dx+b}{cx-a}$

\textbf{Step 3: Case analysis}

\textbf{Case 1: $c=0$}

Function becomes $y = \frac{a}{d}x + \frac{b}{d}$ (linear)

Inverse is $y = \frac{d}{a}x - \frac{b}{a}$

For equality:
\begin{align*}
\frac{a}{d} &= \frac{d}{a} \quad \Rightarrow \quad a^2 = d^2\\
\frac{b}{d} &= -\frac{b}{a} \quad \Rightarrow \quad ab = -bd \quad \Rightarrow \quad b(a+d) = 0
\end{align*}

\textbf{Subcase 1.1:} $b \neq 0$, $d \neq 0$, $a = -d$
\begin{itemize}
    \item 10 choices for $b$ (non-zero: $\pm 1, \pm 2, \ldots, \pm 5$)
    \item 10 choices for $d$ (non-zero: $\pm 1, \pm 2, \ldots, \pm 5$)
    \item $a$ determined by $d$
    \item Count: $10 \times 10 = 100$
\end{itemize}

\textbf{Subcase 1.2:} $b = 0$, $c = 0$, $d \neq 0$, $a^2 = d^2$
\begin{itemize}
    \item 10 choices for $d$ (non-zero)
    \item 2 choices for $a$ (either $d$ or $-d$)
    \item Count: $10 \times 2 = 20$
\end{itemize}

\textbf{Case 1 Total:} $100 + 20 = 120$

\textbf{Case 2: $c \neq 0$}

From $\frac{-dx+b}{cx-a} = \frac{ax+b}{cx+d}$, rewrite in partial fraction form:
\begin{align*}
f(x) &= \frac{a}{c} + \frac{b-\frac{ad}{c}}{cx+d}\\
f^{-1}(x) &= -\frac{d}{c} + \frac{b-\frac{ad}{c}}{cx-a}
\end{align*}

For these to be equal:
\begin{itemize}
    \item Constant terms: $\frac{a}{c} = -\frac{d}{c}$, so $a = -d$
    \item Denominators: $cx+d = cx-a$, so $d = -a$ $\checkmark$
\end{itemize}

Main condition: $a = -d$ with $c \neq 0$

Initial count: $11 \times 11 \times 10 = 1210$
\begin{itemize}
    \item $a \in \{-5, -4, \ldots, 4, 5\}$: 11 choices
    \item $b \in \{-5, -4, \ldots, 4, 5\}$: 11 choices
    \item $c \in \{-5, -4, \ldots, 5\} \setminus \{0\}$: 10 choices
    \item $d = -a$ (determined)
\end{itemize}

\textbf{Exclusions:} When $bc - ad = 0$ with $a = -d$:
\[bc + a^2 = 0 \quad \Rightarrow \quad bc = -a^2\]

This makes the function constant (not symmetric about $y=x$).

Count cases where $bc = -a^2$:

\begin{itemize}
    \item $a=0$: Need $bc=0$ with $c \neq 0$, so $b=0$. Count: 10
    \item $a=\pm 1$: $bc = -1$
    \begin{itemize}
        \item $(b,c) \in \{(1,-1), (-1,1)\}$
        \item Count: $2 \times 2 = 4$
    \end{itemize}
    \item $a=\pm 2$: $bc = -4$
    \begin{itemize}
        \item $(b,c) \in \{(1,-4), (-1,4), (2,-2), (-2,2), (4,-1), (-4,1)\}$
        \item Count: $2 \times 6 = 12$
    \end{itemize}
    \item $a=\pm 3$: $bc = -9$
    \begin{itemize}
        \item $(b,c) \in \{(3,-3), (-3,3)\}$ (since $|b|, |c| \leq 5$)
        \item Count: $2 \times 2 = 4$
    \end{itemize}
    \item $a=\pm 4$: $bc = -16$
    \begin{itemize}
        \item $(b,c) \in \{(4,-4), (-4,4)\}$
        \item Count: $2 \times 2 = 4$
    \end{itemize}
    \item $a=\pm 5$: $bc = -25$
    \begin{itemize}
        \item $(b,c) \in \{(5,-5), (-5,5)\}$
        \item Count: $2 \times 2 = 4$
    \end{itemize}
\end{itemize}

Total exclusions: $10 + 4 + 12 + 4 + 4 + 4 = 38$

\textbf{Case 2 Total:} $1210 - 38 = 1172$

\textbf{Final Answer:}
\[\text{Total} = 120 + 1172 = 1292\]

\subsection{Alternative Methods (Brief Summary)}

\textbf{Method 2: Rotation by $45^\circ$}
\begin{itemize}
    \item Rotate coordinate system by $45^\circ$ counterclockwise
    \item Symmetry about $y=x$ becomes symmetry about $v$-axis
    \item Leads to same condition: $a=-d$
    \item Subtract degenerate cases
    \item Result: $1331 - 28 - 11 - 10 + 10 = 1292$
\end{itemize}

\textbf{Method 3: Asymptote Analysis}
\begin{itemize}
    \item For hyperbola symmetric about $y=x$, asymptotes must intersect on $y=x$
    \item Horizontal asymptote: $y = a/c$
    \item Vertical asymptote: $x = -d/c$
    \item Condition: $a/c = -d/c$, so $a = -d$
    \item Handle $c=0$ case separately
    \item Result: $1292$
\end{itemize}

\textbf{Method 4: Direct Composition}
\begin{itemize}
    \item Compute $f(f(x))$ directly
    \item Set equal to $x$ and expand
    \item Match coefficients: $ac+cd=0$, $d^2-a^2=0$, $ab+bd=0$
    \item Derive conditions and count
    \item Result: $1292$
\end{itemize}

\section{Final Answer}

\begin{center}
\fbox{\Large \textbf{Answer: (B) 1292}}
\end{center}

\vspace{0.5cm}

\textbf{Summary of Solution:}

The condition that the graph of $y = \frac{ax+b}{cx+d}$ is symmetric about $y=x$ is equivalent to the function being its own inverse. Through systematic case analysis:

\begin{itemize}
    \item \textbf{Case 1 ($c=0$):} Linear functions with specific slope/intercept relationships give 120 valid quadruples
    \item \textbf{Case 2 ($c \neq 0$):} General rational functions with $a=-d$ (excluding degenerate cases) give 1172 valid quadruples
\end{itemize}

\textbf{Total: $120 + 1172 = \boxed{1292}$ quadruples}

\end{document}